\section{Numerical illustrations}
\label{sec:numerics}

The experiments are designed as diagnostics for the mechanisms proved above,
not as a search for favorable benchmark instances.  Deterministic Gaussian and
separable low-dimensional log-concave expectations are evaluated analytically
or by fixed Gauss--Hermite quadrature; the nonseparable affine-mode target uses
a fixed scrambled Sobol design.  The stochastic implementations use float64
arithmetic and explicit seeds; baseline/STL algorithm pairs are driven by the
same random-number stream.  The selected figures report the five claims that
matter for the paper; larger parameter sweeps remain in the accompanying
repository.

\subsection{Covariance bootstrap and transport warm start}
\label{sec:numerics-bootstrap}

The covariance-scaling sweep uses diagonal Gaussian targets with initial
covariances as small as \(\lambda_{\min}(C_0)=10^{-10}\); the separate
transport experiment also includes a smooth \(\log\cosh\)-perturbed target,
with \(\lambda_{\min}(C_0)\geq10^{-8}\).  The burn-in endpoint is the first
iterate with \(\lambda_{\min}(C_n)\geq1/(2\beta)\).  Across the Gaussian
scaling sweep, this time is linear in
\(\log(1/(\beta\lambda_{\min}(C_0)))\), as predicted by the two-sided
covariance recursions.  \Cref{fig:num-three-stage} shows one hard KL run with
\(d=20\) and \(\kappa=1000\): covariance calibration precedes rapid energy
decay, which in turn precedes the final decay of the Fisher--Rao mean-block
drift \(\norm{C_n^{1/2}g(a_n)}\).

\begin{figure}[t]
\centering
\begin{subfigure}{0.47\linewidth}
  \centering
  \includegraphics[width=\linewidth]{fig_covboot_three_stage}
  \caption{A hard KL trajectory.}
\end{subfigure}\hfill
\begin{subfigure}{0.51\linewidth}
  \centering
  \includegraphics[width=\linewidth]{fig_covboot_wasserstein_bootstrap}
  \caption{One transport step, then a Fisher--Rao tail.}
\end{subfigure}
\caption{\textbf{Covariance initialization.}
Left: a Gaussian run with \(d=20\), \(\kappa=1000\), and
\(\lambda_{\min}(C_0)=10^{-10}\) separates covariance burn-in, global energy
localization, and the final local regime.  Right: one
Bures--Wasserstein bootstrap with \(c=1/2\), hence transport step
\(\eta=c/\beta\), lifts the minimum covariance eigenvalue to the curvature
scale independently of the initial underdispersion; the subsequent
Fisher--Rao warm-up is essentially flat in \(\lambda_{\min}(C_0)\).}
\label{fig:num-three-stage}
\end{figure}

The transport experiment makes a deliberately narrow point.  A single
Bures--Wasserstein step removes the small-covariance cost and then hands the
trajectory back to a Fisher--Rao scheme.  It neither changes the
\(\kappastar\)-scale spiral obstruction nor implements a blended WFR flow.

\subsection{Riemannian versus KL discretization}
\label{sec:numerics-discretization}

The deterministic comparison uses two-dimensional Gaussian and smooth
log-concave targets, together with a convex quartic target whose Hessian is
unbounded and hence is not globally smooth.  Both schemes use the same explicit
mean update and evaluate the same population score and Hessian, analytically or
by deterministic quadrature.  The Riemannian covariance update applies a
matrix-exponential retraction; the KL update solves a rational precision
resolvent.  We report iteration count and elapsed flow time \(nh\)
separately; the latter is a discretization diagnostic, not wall-clock timing.

\begin{figure}[t]
\centering
\begin{subfigure}{0.58\linewidth}
  \centering
  \includegraphics[width=\linewidth]{fig_convergence_speed_smooth_logconcave}
  \caption{Matched-step trajectories on a smooth target.}
\end{subfigure}\hfill
\begin{subfigure}{0.39\linewidth}
  \centering
  \includegraphics[width=\linewidth]{fig_iterations_vs_stepsize}
  \caption{Iterations to tolerance.}
\end{subfigure}
\caption{\textbf{Deterministic discretizations.}
At matched stable stepsizes, the Riemannian and KL trajectories have comparable
first-order behavior.  The Riemannian map often contracts faster per step in
the tested range, whereas the KL covariance resolvent remains well defined for
all stepsizes on log-concave targets and avoids a matrix exponential.  Both
methods can still become unstable through their shared explicit mean update.
The common theorem-safe stepsize is sufficient and conservative in these
low-dimensional examples.}
\label{fig:num-discretization}
\end{figure}

The descriptor \emph{more robust} in this comparison refers specifically to
the covariance map.  On a log-concave target,
\(C_n^{-1}+hA_n\succ0\) for every \(h>0\), so the KL covariance update cannot
leave the positive-definite cone.  The full algorithm may nevertheless fail
at a large step because the mean update contains the factor \(hC_n\).

\subsection{STL variance and constant-step floors}
\label{sec:numerics-stl}

The stochastic grid uses Gaussian targets with
\(d\in\{2,8,32\}\), \(\kappa\in\{10,100,10^4\}\), and
smooth \(\log\cosh\) targets with
\(d\in\{2,8,32\}\), \(\kappa\in\{10,100\}\).  Each cell uses 20 seeds.
Baseline and STL runs share the same samples, so the comparison is paired.
The energy shown in the plots is evaluated independently of the update sample.

\begin{figure}[t]
\centering
\begin{subfigure}{0.48\linewidth}
  \centering
  \includegraphics[width=\linewidth]{fig_stl_variance_by_distance_to_optimum}
  \caption{Fisher--Rao variance ratio versus mean displacement.}
\end{subfigure}\hfill
\begin{subfigure}{0.49\linewidth}
  \centering
  \includegraphics[width=\linewidth]{fig_stl_batchsize_effect}
  \caption{Tail energy versus batch size.}
\end{subfigure}
\caption{\textbf{Price/Hessian STL diagnostics.}
For a non-Gaussian \(\log\cosh\) target, the displayed mean-block Fisher--Rao
variance ratio \(\Var_{\FR}(b_{\STL})/\Var_{\FR}(b_{\rm base})\) is positive at
the optimizer and increases with the mean displacement in every plotted
configuration (left).  Mini-batching reduces the median tail energy gap
approximately as \(1/B\), while STL stays below its paired baseline in the
recorded grid (right).  For a Gaussian target with matched covariance, the
mean estimator and the constant Hessian are deterministic pathwise; this
special calibration is not extrapolated to general targets.}
\label{fig:num-stl}
\end{figure}

These plots support the variance decomposition of
\Cref{sec:stochastic}.  They do not test the high-probability local
non-exit theorem, whose horizon-dependent batch requirement is a worst-case
containment statement rather than a prediction for the median trajectory.

\subsection{Trace and traceless affine modes}
\label{sec:numerics-affine}

The affine-metric experiment varies
\(\tau\in\{-\omega/(2d),0,\omega/(2d)\}\) for several values of \(\omega\).
Five initializations isolate mean, volume-high, volume-low, traceless shape,
and mixed perturbations.  The Gaussian target uses exact expectations.  The
non-Gaussian target is strongly log-concave and uses one common fixed,
scrambled Sobol quasi--Monte Carlo design across all metric choices.

\begin{figure}[t]
\centering
\begin{subfigure}{0.49\linewidth}
  \centering
  \includegraphics[width=\linewidth]{fig_ot_gaussian_tau_speedup}
  \caption{Gaussian target.}
\end{subfigure}\hfill
\begin{subfigure}{0.49\linewidth}
  \centering
  \includegraphics[width=\linewidth]{fig_ot_logconcave_tau_speedup}
  \caption{Strongly log-concave target.}
\end{subfigure}
\caption{\textbf{Affine-invariant mode preconditioning.}
The panels show the time-to-tolerance ratio \(T(\tau)/T(0)\).  For
\(\tau=-\omega/(2d)\), the trace rate doubles and the volume-dominated runs
take about half as long.  For \(\tau=\omega/(2d)\), their time scale grows by
about \(3/2\).  Mean- and traceless-dominated runs remain near one.  The
Gaussian panel calibrates the exact modal calculation; the non-Gaussian panel
is an illustration, not a general convergence theorem for the full metric
family.}
\label{fig:num-affine-modes}
\end{figure}

\subsection{The non-log-concave boundary}
\label{sec:numerics-nonconvex}

For \eqref{eq:nonconvex-potential}, all expectations are deterministic
Gauss--Hermite quadratures and all trajectories are centered.  The first two
panels visualize the analytic mechanisms of
\Cref{prop:nonconvex-riemannian-cascade,prop:nonconvex-kl-pole}.  The last
panel evaluates the covariance part of the split stationarity measure in
\eqref{eq:nonconvex-split-stationarity}; the mean term is zero by symmetry.

\begin{figure}[t]
\centering
\begin{subfigure}{0.32\linewidth}
  \centering
  \includegraphics[width=\linewidth]{fig_nonconvex_riemannian_cascade}
  \caption{Riemannian cascade.}
\end{subfigure}\hfill
\begin{subfigure}{0.32\linewidth}
  \centering
  \includegraphics[width=\linewidth]{fig_nonconvex_kl_pole}
  \caption{KL pole.}
\end{subfigure}\hfill
\begin{subfigure}{0.32\linewidth}
  \centering
  \includegraphics[width=\linewidth]{fig_nonconvex_clipped_kl_stationarity}
  \caption{Clipped-KL stationarity.}
\end{subfigure}
\caption{\textbf{Beyond log-concavity.}
Although \(\abs{V_R''}\leq1\), negative curvature drives a covariance cascade
in the Riemannian map and a rational pole in the unprojected KL map.  With
\(C_n\in[\lambda_-,\lambda_+]\) and a theorem-safe step, the running minimum
of the centered covariance KL stays below the energy-drop envelope from
\Cref{thm:nonconvex-clipped-stationarity}.  After the upper constraint has
bound, the iterate remains pinned there and the next displacement is zero;
this is constrained stationarity, not a small unconstrained Fisher--Rao
gradient.}
\label{fig:num-nonconvex}
\end{figure}

\subsection{Reproducibility}
\label{sec:numerics-reproducibility}

The repository contains the experiment scripts, the raw outputs used by the
selected panels, and a deterministic asset-building script.  The selected
submission PDFs are byte-identical to the corresponding report assets.
The selected panels can be rebuilt directly from their surviving CSV inputs,
but the covariance-bootstrap and STL group-level builder entry points do not
run end to end from the current snapshot because some supplementary trajectory
and run-metadata files are absent.
\Cref{app:experiments} records the target formulas, grids, seeds, quadrature
rules, stopping conditions, exact data-to-figure lineage, and the provenance
files that are absent from the current repository snapshot.  No general WFR
output is used in this paper.
